\documentclass[10pt,letterpaper]{article}
\usepackage[utf8]{inputenc}
\usepackage[vietnamese]{babel}
\usepackage{amsmath,amssymb,amsfonts}
\usepackage{graphicx}
\usepackage{xcolor}
\usepackage{tikz}
\usepackage{enumitem}
\usepackage{fancyhdr}
\usepackage{calc}
\usepackage{newtxtext,newtxmath}

\usepackage[top=1.8cm, bottom=2.0cm, left=1.6cm, right=1.6cm, headheight=14pt]{geometry}

% Định nghĩa màu sắc theo phong cách sách Stewart Calculus
\definecolor{stewartcyan}{RGB}{0, 136, 170}
\definecolor{stewartred}{RGB}{185, 30, 50}

% Cấu hình header & footer
\pagestyle{fancy}
\fancyhf{}
\renewcommand{\headrulewidth}{0pt}
\fancyhead[R]{\fontsize{9pt}{11pt}\selectfont \textbf{\textsf{MỤC 4.3}}\quad \textsf{Đạo hàm cho ta biết điều gì về hình dạng của đồ thị}\qquad \textbf{305}}
\fancyfoot[C]{\fontsize{6.5pt}{8pt}\selectfont \color{gray!80}%
Copyright 2021 Cengage Learning. All Rights Reserved. May not be copied, scanned, or duplicated, in whole or in part. WCN 02-200-203\\% 
Editorial review has deemed that any suppressed content does not materially affect the overall learning experience. Cengage Learning reserves the right to remove additional content at any time if subsequent rights restrictions require it.}

\setlength{\parindent}{0pt}
\setlength{\parskip}{0pt}

\begin{document}

% Đoạn nối tiếp từ trang trước nằm ở phần cột chính (x >= 180pt)
\noindent
\hspace*{0.26\textwidth}%
\begin{minipage}{0.74\textwidth}
    phác thảo bằng cách kết hợp thông tin liên quan đến tính lõm và điểm uốn. Trong Hình 14(c), chúng ta kiểm tra lại kết quả của mình bằng máy tính.
\end{minipage}

\vspace{0.7em}

% HÌNH 14 (FIGURE 14)
\noindent
\begin{minipage}{\textwidth}
    \centering
    \includegraphics[width=\linewidth]{images/p0340_fig1.png}
\end{minipage}

\vspace{0.3em}
\noindent
\textbf{\textsf{HÌNH 14}} \hfill {\color{stewartcyan}\rule{1.1ex}{1.1ex}}

\vspace{1.0em}

% BANNER TIÊU ĐỀ: 4.3 BÀI TẬP (4.3 EXERCISES)
\noindent
\begin{tikzpicture}
    \pgfmathsetmacro{\W}{\textwidth/1cm}
    % Hai vạch kẻ bên trái
    \draw[line width=1.1pt, stewartcyan] (0, 0.26) -- (1.5, 0.26);
    \draw[line width=1.1pt, stewartcyan] (0, 0.12) -- (1.5, 0.12);
    
    % Tiêu đề chính giữa
    \node[anchor=west, inner sep=3pt] (t) at (1.55, 0.19) {
        \textbf{\textsf{\textcolor{stewartcyan}{\LARGE 4.3}}}\;\;%
        {\color{stewartcyan!80}\rule[-0.25ex]{1.3pt}{1.35em}}\;\;%
        \textbf{\textsf{\LARGE Bài tập}}
    };
    
    % Ba vạch kẻ bên phải kéo dài đến hết lề
    \path (t.east);
    \pgfgetlastxy{\tx}{\ty}
    \pgfmathsetmacro{\xR}{\tx/28.4527 + 0.15}
    \draw[line width=1.1pt, stewartcyan] (\xR, 0.33) -- (\W, 0.33);
    \draw[line width=1.1pt, stewartcyan] (\xR, 0.19) -- (\W, 0.19);
    \draw[line width=1.1pt, stewartcyan] (\xR, 0.05) -- (\W, 0.05);
\end{tikzpicture}

\vspace{0.8em}

% BỐ CỤC 2 CỘT CHO BÀI TẬP
\noindent
\begin{minipage}[t]{0.485\textwidth}
    % CỘT TRÁI
    \noindent
    \textbf{\textcolor{stewartcyan}{1--2}}\;\; Sử dụng đồ thị đã cho của hàm số $f$ để tìm các thông tin sau:
    \begin{enumerate}[label=\textbf{(\alph*)}, itemsep=1pt, topsep=2pt, parsep=0pt, leftmargin=1.6em]
        \item Các khoảng mở mà trên đó $f$ đồng biến.
        \item Các khoảng mở mà trên đó $f$ nghịch biến.
        \item Các khoảng mở mà trên đó $f$ lõm lên trên.
        \item Các khoảng mở mà trên đó $f$ lõm xuống dưới.
        \item Tọa độ các điểm uốn.
    \end{enumerate}

    \vspace{0.4em}
    \noindent
    \begin{minipage}[t]{0.48\linewidth}
        \textbf{1.}\vspace{-1.2em}
        \begin{center}
            \includegraphics[width=0.95\linewidth]{images/p0340_fig2.png}
        \end{center}
    \end{minipage}\hfill
    \begin{minipage}[t]{0.48\linewidth}
        \textbf{2.}\vspace{-1.2em}
        \begin{center}
            \includegraphics[width=0.95\linewidth]{images/p0340_fig3.png}
        \end{center}
    \end{minipage}

    \vspace{0.8em}
    \noindent
    \textbf{3.}\;\; Giả sử bạn được cho công thức của một hàm số $f$.
    \begin{enumerate}[label=\textbf{(\alph*)}, itemsep=1pt, topsep=2pt, parsep=0pt, leftmargin=1.6em]
        \item Làm thế nào để bạn xác định được $f$ đồng biến hay nghịch biến trên những khoảng nào?
        \item Làm thế nào để bạn xác định được đồ thị của $f$ lõm lên trên hay lõm xuống dưới?
        \item Làm thế nào để bạn xác định vị trí các điểm uốn?
    \end{enumerate}

    \vspace{0.7em}
    \noindent
    \textbf{4.}\;\; \textbf{(a)}\; Phát biểu Kiểm định đạo hàm cấp một.
    \begin{enumerate}[label=\textbf{(\alph*)}, itemsep=1pt, topsep=2pt, parsep=0pt, leftmargin=1.6em, start=2]
        \item Phát biểu Kiểm định đạo hàm cấp hai. Trong những trường hợp nào thì kiểm định này không đưa ra kết luận? Bạn sẽ làm gì nếu nó không hiệu quả?
    \end{enumerate}

    \vspace{0.7em}
    \noindent
    \textbf{\textcolor{stewartcyan}{5--6}}\;\; Đồ thị của đạo hàm $f'$ của một hàm số $f$ được cho như hình vẽ.
    \begin{enumerate}[label=\textbf{(\alph*)}, itemsep=1pt, topsep=2pt, parsep=0pt, leftmargin=1.6em]
        \item Trên những khoảng nào thì $f$ đồng biến? Nghịch biến?
        \item Tại những giá trị nào của $x$ thì $f$ có cực đại địa phương? Cực tiểu địa phương?
    \end{enumerate}

    \vspace{0.2em}
    \noindent
    \textbf{5.}\vspace{-1.2em}
    \begin{center}
        \includegraphics[width=0.76\linewidth]{images/p0340_fig6.png}
    \end{center}
\end{minipage}%
\hfill
\begin{minipage}[t]{0.485\textwidth}
    % CỘT PHẢI
    \noindent
    \textbf{6.}\vspace{-1.2em}
    \begin{center}
        \includegraphics[width=0.76\linewidth]{images/p0340_fig4.png}
    \end{center}

    \vspace{0.8em}
    \noindent
    \textbf{7.}\;\; Trong mỗi trường hợp, hãy nêu hoành độ $x$ của các điểm uốn của $f$. Nêu lý do cho câu trả lời của bạn.
    \begin{enumerate}[label=\textbf{(\alph*)}, itemsep=1pt, topsep=2pt, parsep=0pt, leftmargin=1.6em]
        \item Đường cong là đồ thị của hàm số $f$.
        \item Đường cong là đồ thị của đạo hàm $f'$.
        \item Đường cong là đồ thị của đạo hàm cấp hai $f''$.
    \end{enumerate}

    \vspace{-0.3em}
    \begin{center}
        \includegraphics[width=0.82\linewidth]{images/p0340_fig5.png}
    \end{center}

    \vspace{0.3em}
    \noindent
    \textbf{8.}\;\; Đồ thị của đạo hàm cấp một $f'$ của một hàm số $f$ được cho như hình vẽ.
    \begin{enumerate}[label=\textbf{(\alph*)}, itemsep=1pt, topsep=2pt, parsep=0pt, leftmargin=1.6em]
        \item Trên những khoảng nào thì $f$ đồng biến? Giải thích.
        \item Tại những giá trị nào của $x$ thì $f$ đạt cực đại địa phương hoặc cực tiểu địa phương? Giải thích.
        \item Trên những khoảng nào thì $f$ lõm lên trên hoặc lõm xuống dưới? Giải thích.
        \item Hoành độ $x$ của các điểm uốn của $f$ là những giá trị nào? Tại sao?
    \end{enumerate}

    \vspace{-0.3em}
    \begin{center}
        \includegraphics[width=0.82\linewidth]{images/p0340_fig7.png}
    \end{center}
\end{minipage}

\end{document}